\documentclass[12pt,letterpaper]{article}
\usepackage[T1]{fontenc}
\usepackage[letterpaper,margin=1in]{geometry}
\usepackage{amsmath,amssymb,amsthm,array,booktabs,microtype,longtable}
\IfFileExists{newtxtext.sty}{\usepackage{newtxtext,newtxmath}}{\usepackage{mathptmx}}
\usepackage[hidelinks]{hyperref}
\title{Revision Log: Representations of finite groups on the Riemann--Roch space of a toric variety\\\large Draft 4, corrected and expanded version}
\author{David Joyner (original manuscript); revision assistance by ChatGPT}
\date{November 22, 2026}
\begin{document}
\maketitle
\section{Provenance and scope}
The source was the original LaTeX manuscript \texttt{rr\_toric.tex}, dated June 23, 2002, supplied in the conversation, together with a referee report by Carbon (Claude), supplied as a text file. This revision preserves the paper's title, problem, induced action on the Riemann--Roch space, and its original two-dimensional toric example. It corrects mathematical errors and adds a graph-theoretic treatment. It is not represented as a peer-reviewed or fully novel research contribution.

\section{Responses to Carbon's report}
\begin{longtable}{@{}>{\raggedright\arraybackslash}p{0.32\linewidth}>{\raggedright\arraybackslash}p{0.60\linewidth}@{}}
\toprule
Issue & Resolution in Draft 2\\\midrule
\endhead
Orbit representations falsely called irreducible & Replaced by the induced permutation module formula and a separate irreducible-character multiplicity formula.\\
Theorem too broadly assumes arbitrary automorphisms & Restricted the permutation theorem to finite groups of fan automorphisms preserving a torus-invariant divisor. Added weighted fixed-point trace for finite subgroups of the torus normalizer.\\
Smoothness versus singular example & Stated the theorem for complete normal toric varieties and arbitrary torus-invariant Weil divisors, using the reflexive sheaf $\mathcal O_X(D)$.\\
Incorrect affine chart equations & Replaced all three with cubic $A_2$ hypersurface equations $u^3=vw$; clarified the invariant and exchanged charts.\\
Dual lattice and orbit notation & Distinguished $N$ from its dual $M=N^\vee$, and distinguished the torus orbit $O(\rho)$ from its closure $D_\rho=V(\rho)$.\\
Proof incorrectly invokes Maschke & Supplied a direct orbit-basis proof. Maschke is used only for subsequent irreducible decomposition over characteristic zero.\\
Trace theorem insufficiently explained & Proved the fixed-lattice-point character formula and irreducible multiplicity formula using character inner products.\\
Misprints, unused symbols and outdated bibliography & Replaced inconsistent $L^\perp$ notation, removed unused inner product and unused quantifier, corrected mathematical typos, and supplied updated bibliography with stable keys.\\
Concern about novelty & Literature review acknowledges established equivariant Ehrhart theory and graph-zonotope work. The revision does not characterize standard formulas as new.\\
\bottomrule
\end{longtable}

\section{Original mathematical results retained or corrected}
The 2002 note contained the basic idea of decomposing the character basis by finite-group orbits and calculating the trace from fixed lattice points. That insight is preserved, but the asserted irreducibility of each orbit module is false. The original fan with rays $(2,-1),(-1,2),(-1,-1)$ is valid and complete, though singular. The original divisor $D=6D_1+6D_2$ and triangle with vertices $(-6,-6),(-2,2),(2,-2)$ are retained. The numbers of fixed and two-element orbits, $7$ and $12$, are correct.

\section{Results and expositions introduced in the correction}
\begin{enumerate}
\item The decomposition $L(D)=\bigoplus\operatorname{Ind}_{G_u}^{G}\mathbf1$ is stated and proved correctly, followed by an irreducible multiplicity formula. These are standard consequences of toric monomial bases and character theory, not claims of first discovery.
\item For all $q\geq0$, the original surface yields $\dim L(qD)=24q^2+6q+1$, involution trace $6q+1$, and multiplicities $12q^2+6q+1$ and $12q^2$.
\item Lattice-point graphs with $G$-invariant symmetric step sets yield commuting adjacency operators and a twisted-walk trace formula. These consequences are proved directly and are presented as elementary graph-representation constructions.
\item The $31$-vertex graph on the original triangle has $85$ edges and twisted adjacency trace $10$.
\item A graphic-zonotope section relates toric section dimensions to forest enumerators. This is standard literature, cited explicitly.
\item For $K_3$, character and irreducible multiplicity formulas for the hexagon and its dilates are derived.
\item For the hypersimplex $\Delta(2,n)$, the $S_n$ edge-permutation representation is decomposed using the Johnson graph $J(n,2)$. The spectrum and irreducibility are established with an unsigned incidence-matrix identity and a three-dimensional commutant argument. These are classical facts organized as a toric example.
\item The torus-normalizer trace is treated separately, to prevent confusing permutation with monomial actions.
\end{enumerate}

\section{Numbering concordance}
\begin{center}
\begin{tabular}{@{}p{0.46\linewidth}p{0.45\linewidth}@{}}
\toprule
2002 manuscript & Draft 2 manuscript\\\midrule
Theorem 1 (incorrect irreducibility) & Theorem 4.1 (corrected orbit theorem) and Corollary 4.2\\
Theorem 2 (trace for fan automorphisms) & Corollary 4.2 (character formula)\\
Section 2 (toric surface example) & Section 5 (repaired surface example)\\
Inverse question after Theorem 1 & Section 10 (summary and open questions)\\
No original counterparts & Sections 6--9 (graph theory and monomial extension)\\
\bottomrule
\end{tabular}
\end{center}

\section{Verification and source checking}
\begin{itemize}
\item \textbf{Level 1: independent mathematical derivation.} The orbit-module and trace formulas, lattice area/boundary counts by Pick's theorem, $C_2$ multiplicities, commuting adjacency and twisted-walk identities, K3 fixed-point counts, weighted torus-normalizer trace, and Johnson incidence-matrix proof were derived from the specified definitions.
\item \textbf{Level 1: exact computational checks.} The accompanying Python script (\path{verify_examples.py}) enumerates $P_{qD}\cap\mathbb Z^2$ and K3 hexagon lattice points for $0\leq q\leq12$, checks the graph's $85$ edges, connectedness, five swapped edges, and twisted trace $10$, and checks Johnson incidence-matrix identities for $4\leq n\leq12$.
\item \textbf{Level 2: consulted sources.} The public full text of Elia--Kim--Supina (2024) was consulted for its account of equivariant Ehrhart theory, graphic zonotopes, and hypersimplices. Bibliographic metadata and abstracts were checked for Stapledon (2011), Joyner--Traves (2002), Brouwer--Haemers (2012), Oda (1988), and James (1978). A mathematical exposition of the forest-counting identity was also consulted.
\item \textbf{Level 3: classical references not fully checked.} The full texts of Fulton (1993), Serre (1977), James (1978), and Brouwer--Haemers (2012) were not systematically checked line by line. The use of classical toric and representation facts is accompanied by self-contained arguments where feasible. Thus no source-checking claim beyond the material actually consulted is made.
\item \textbf{Typesetting.} Draft 2 compiles with \texttt{pdflatex} using 12pt Times-style text and math, US Letter paper, and one-inch margins. Equation references and bibliography keys were checked by recompilation.
\end{itemize}

\section{AI usage and limitations}
Carbon (Claude) supplied an initial referee report, which motivated the correction of the false irreducibility claim and erroneous affine equations. ChatGPT (OpenAI GPT-6) conducted the subsequent revision, developed the graph sections, checked mathematical derivations, wrote the exact verification script, checked bibliographic sources, and compiled the LaTeX. These AI activities are not an independent certification of the results. Responsibility for mathematical verification and any publication or submission remains with the human author.

\section{Draft 3: expository separation}
All historical criticism of the 2002 manuscript and the complete response to Carbon's
referee report remain in the preceding sections of this log. Draft 3 presents the
mathematical material without references to errors in earlier drafts, retaining
the same results, examples, proofs, and bibliographic citations as Draft 2.

The introduction now poses the natural question of how orbit modules on lattice
points determine irreducible character multiplicities. The surface section
introduces the singular fan and its explicit affine charts as independent
calculations, without contrasting them with earlier equations. The graph-theory
sections proceed directly from commuting operators to the Johnson-graph
example. The torus-normalizer section asks how the fixed-point character formula
changes when an action includes torus translations, leading to weighted traces.
Historical reasons for distinguishing transitive permutation modules from
irreducibles, and fan automorphisms from general toric automorphisms, remain
recorded above rather than in the body of the paper.

\noindent\textbf{Status and verification.} This is an editorial revision.
No theorem, numerical formula, proof, or mathematical example was intentionally
changed. The exact verification script was rerun and the manuscript and log
were recompiled with \texttt{pdflatex}. AI review is not an independent
certification of correctness.


\section{Draft 4: response to Claude's November 22, 2026 report}
This section records the changes made after Claude (Opus 5.5) reviewed Draft 3.
The previous sections remain a historical record of earlier work. Draft 4
continues to present the mathematical questions naturally, with discussion
of past mistakes confined to this log.

\subsection*{Required mathematical repairs (independently checked)}
\begin{enumerate}
\item \textbf{Inverse problem, Section 10 (false as phrased in Draft 3).}
The pair $(1,5)$, $(2,3)$ on $\mathbb P^1\times\mathbb P^1$ gives equal
section dimensions only for $q=1$: at $q=2$ the values are $33$ and $35$.
It has been replaced by the nonequivalent bidegrees $(1,2)$ and $(2,1)$,
for which both dimension sequences are $(q+1)(2q+1)$. These divisors
are exchanged by a factor-swapping fan automorphism. The resulting natural
open question now asks for recovery modulo appropriate normalizing
fan automorphisms. Verification level 1 (direct calculation).
\item \textbf{Affine symmetries and linearization, Sections 4, 7, 8.}
The original theorem assumes a linear action on the character lattice
preserving the divisor coefficients. Polytopes in an affine lattice,
such as the hypersimplex, need the additional construction now stated
as Theorem 4.3. Its proof uses the monomial graded section ring:
in degree $q$ an affine map $u\mapsto A_gu+b_g$ acts by
$u\mapsto A_gu+qb_g$. This produces the permutation basis for a
specified monomial linearization. Other linearizations of the same
line bundle and base action may twist the representation by a group
character. Verification level 1 (direct proof).
\end{enumerate}

\subsection*{Other accepted corrections}
\begin{enumerate}
\item The third affine chart relation has been corrected to
$z_1^3=z_2z_3$ in the specified monomial order, and the exponent
relation is shown. The $K_3$ zonotope is translated by
$-(1,1,1)$, not by $(1,1,1)$.
\item The Johnson incidence-matrix proof now explains the standard
eigenspace and injectivity of $R^{\mathsf T}$, and identifies
$\ker R$ with the Specht module $S^{(n-2,2)}$. The $n\geq4$
hypothesis is tied to the disjoint-pair orbital, not the rank claim.
James (1978) is cited for the Specht module description.
\item The invariant Weil-divisor basis formula now cites
Cox--Little--Schenck, Proposition 4.3.3.
Stanley (1991) is cited for graphic-zonotope enumeration.
Delsarte (1973) is cited for the classical Johnson association scheme.
The forest-counting identity is stated for arbitrary simple graphs;
connectedness is retained only where full dimension of the graphic
zonotope is required.
\item The source now distinguishes $\rho$ (fan rays) from $\pi_D$
(the representation), defines the adjacency eigenvalue set,
removes a redundant $C_2$ example, repairs several glossary
definitions, and gives a concrete representation count in the introduction.
\end{enumerate}

\subsection*{Accepted strengthenings}
\begin{enumerate}
\item A non-Cartier divisor $D'=2D_1+2D_2$, with $3D'=D$, now provides
the exact period-three quasipolynomial
\[
\dim L(qD')=\frac{8q^2+6q+1}{3}+\frac23\mathbf1_{3\mid q},
\qquad \operatorname{tr}(g\mid L(qD'))=2q+1.
\]
The paper includes a direct parity-and-floor-sum proof; exact enumeration
was independently checked for $0\leq q\leq9$.
\item The 31-vertex graph's even and odd adjacency characteristic
polynomials were independently computed over the rationals.
Each is squarefree, and their greatest common divisor is one.
Thus all 31 eigenvalues are simple. This is an exact computational
observation, not a new general spectral theorem.
\item The root step set $B=\{e_i-e_j:i\ne j\}$ identifies $J(n,2)$
as an affine lattice-point graph. A remark describes the classical
$k$-subset and Johnson-scheme extension for $\Delta(k,n)$.
No originality is claimed for these decompositions.
\end{enumerate}

\subsection*{Verification limits and editorial decisions}
The displayed correction to the inverse example and the non-Cartier
formula were checked by direct calculation. The characteristic-polynomial
claim was checked independently using exact symbolic integer arithmetic,
without floating-point eigenvalue approximations.
The 14-page manuscript compiles with \texttt{pdflatex} (three passes),
12pt Times-style text and math, US Letter, no undefined references
and no overfull boxes. The verification script in the repository
supports repeatable checks. Bibliographic metadata for Stanley (1991)
were checked against his publication list, and Delsarte (1973) against
catalog records. The full texts of James (1978), Cox--Little--Schenck (2011),
and Delsarte (1973) were not independently examined line by line, so claims
about their detailed contents are not represented as level-2 source checks.
No material recommendation from Claude's latest report was rejected.
The text retains the author's 2002 mathematical framework, while the
graph-theoretic and affine-linearization developments are from the 2026
revisions and should not be attributed to the original manuscript.

\subsection*{AI assistance}
Claude (Opus 5.5) supplied a Draft 3 referee report, including the
scope gap, error diagnosis and proposed strengthenings.
ChatGPT (OpenAI, GPT-6) independently recalculated the inverse example,
the non-Cartier surface count and the adjacency characteristic-polynomial
properties; prepared the Draft 4 revision, bibliography, proofs,
scripts and repository updates. AI review does not constitute independent
certification of mathematical correctness.

\section{Draft 4: explanatory glossary expansion (November 22, 2026)}
This editorial pass adds two explanatory sentences to each of the twelve
glossary entries. The additions supply concrete formulas, examples, and
connections to the manuscript: monomial characters, fixed-point traces,
graphic-zonotope forests, hypersimplex subsets, coset modules, and
twisted walks. The original definitions and their terminology are retained.
No theorem, proof, mathematical formula outside the glossary, or bibliography
entry was altered in this pass.

\noindent\textbf{Verification.} The revised manuscript was compiled locally
with \texttt{pdflatex} in two passes. The compilation reported no warnings
or overfull boxes. This revision is expository only and does not upgrade
any proof or verification level.

\section{Suggested next review}
A human referee should pay particular attention to the natural $S_n$-linearization on the hypersimplex line bundle, to the scope of the zonotope lattice-count identity in the intrinsic affine lattice, and to prior art before attributing originality to any graph-spectral generalization.
\end{document}
